
In this section, we will review the ABC framework and highlight its limitations within the hierarchical model framework. Subsequently, we will introduce our permutation-based method, permABC, which aims to address these limitations.

\subsection{Introduction to ABC}


The standard rejection ABC algorithm, often referred to as \emph{Vanilla ABC} (see Algorithm~\ref{alg:abc_vanilla}), proceeds by drawing parameters from the prior distribution, $\param \sim \pi(\cdot)$, and then simulating synthetic data from the likelihood, $\simdata \sim f(\cdot \mid \param)$. The pair $(\param, \simdata)$ is accepted if the summary statistics of the simulated data are sufficiently close to those of the observed data, that is, if $d\{\eta(\obsdata), \eta(\simdata)\} \leq \varepsilon$ for some tolerance level $\varepsilon > 0$ and some summary statistic $\eta$. This procedure yields samples from the approximate joint distribution $\pi_{\varepsilon}\{\param,\simdata \mid \eta(\obsdata)\}$; we marginalize out $\simdata$ to get the ABC posterior

\begin{equation}
    \label{eqabc}
    \begin{split}
        \pi_{\varepsilon}\{\param \mid \eta(\obsdata)\} &\propto
        \int_{\dataspace^\Clust} \pi_{\varepsilon}\{\param, \simdata \mid \eta(\obsdata)\} \, d\simdata \\
        &\propto \pi(\param) \int_{A_{\obsdata,\varepsilon}} f(\simdata \mid \param) \, d\simdata,
    \end{split}
\end{equation}
where $A_{\obsdata,\varepsilon} = \{\simdata \in \dataspace^\Clust \mid d\{\eta(\obsdata), \eta(\simdata)\} \leq \varepsilon\}$ denotes the acceptance region.

Combining an informative summary statistic $\eta$ with a small tolerance $\varepsilon$ typically leads to a good approximation of the true posterior distribution, so that $\pi_{\varepsilon}\{\param \mid \eta(\obsdata)\} \approx \pi(\param \mid \obsdata)$. This approximate posterior—often called a \emph{pseudo-posterior}—is obtained by replacing the true likelihood $f(\obsdata \mid \param)$ with the integral of the likelihood over the neighborhood $A_{\obsdata,\varepsilon}$ of the observed data, effectively using it as a surrogate for the intractable likelihood.

Throughout this work, we assume for clarity that the data are compared directly, i.e., $\eta = \mathrm{Id}$. Under this assumption, the acceptance region simplifies to the $\varepsilon$-ball $B_{\varepsilon}(\obsdata)$ centered at $\obsdata$ with respect to the distance $d$. Nonetheless, the proposed method naturally extends to the use of low-dimensional summary statistics.


In our numerical examples, we take $\dataspace = \mathbb R^n$ and use the distance $d(\cdot,\cdot)$ between two datasets in \( \dataspace^K \), each consisting of \( K \times n \) observations, defined as follows: 
\begin{equation}
    \label{eqdistance}
    d(\obsdata, \simdata)^2 = {\sum_{k=1}^K w_k^2 \| \obsdata_k - \simdata_k \|_2^2},
\end{equation}
where \( \| \cdot \|_2 \) denotes the Euclidean norm, and the weights \( w_k \) are used to control the relative influence of each compartment in the distance computation. The weight vector \( w_{1:K} \) can either be manually specified by the user or selected adaptively to prevent any single compartment from dominating the distance, as suggested by \citet{prangle2017adapting}.


This naive ABC algorithm faces several limitations, including poor scalability in the parameter space and an exponentially increasing computational cost due to the declining acceptance rate as \( \varepsilon \) decreases. In the context of the hierarchical model of Figure \ref{fig:hier}, the acceptance rate further diminishes as the number of compartments \( K \) increases. Simulating a dataset \( \simdata \) from the prior predictive distribution that matches all compartments of \( \obsdata \) within a small tolerance \( \varepsilon \) can be particularly challenging. 



\begin{algorithm}
\label{alg:abc_vanilla}
\caption{Vanilla Approximate Bayesian Computation}
\textbf{Input: }observed dataset $\obsdata$, number of iterations $N$ , threshold $\eps>0$, summary statistic $\eta$.\\
\textbf{Output: }a sample $(\param^{(1)}, \dots, \param^{(N)})$ from $\pi_\eps\{\cdot\mid\eta(\obsdata)\}$.\\

\BlankLine

\For{$i = 1,\dots, N$}{
    \Repeat{$d\{\eta(\simdata^{(i)}),\eta(\obsdata)\})\leq \eps$}{
        $\param^{(i)} \sim \pi(\cdot)$\;
        $\simdata^{(i)}\sim f(\cdot\mid\param^{(i)})$\;
        }
    }
\end{algorithm}


\subsection{Motivations of permABC}

A common approach in ABC to increase the acceptance rate is to exploit the exchangeability of the compartments by using a permutation-invariant summary statistic, such as the sorted vector. For instance, when we have a one-dimensional observation (or statistic) per compartment, $\obsdata \in \mathbb{R}^K$, it is natural to use the ordering $\eta(\obsdata) = \text{sort}(\obsdata)$ as a summary. \old{This reduces the effective complexity from $K!$ permutations to a canonical representation, thereby simplifying the comparison between datasets.}

Moreover, in the univariate case, sorting is known to minimize the distance over all permutations of simulated data. Specifically, $d\{\text{sort}(\obsdata), \text{sort}(\simdata)\} = \min_{\sigma \in \mathcal{S}_K} d(\obsdata, \simdata_{\sigma})$, where $\mathcal{S}_K$ denotes the set of all permutations of $\{1,\dots,K\}$. This property underlies several recent developments in ABC, including those based on optimal transport and Wasserstein distances \citep{bernton2019approximate}, where sorting plays a central role in defining distances between empirical distributions.

\new{This perspective is closely related in spirit to Wasserstein ABC (WABC) \citep{bernton2019approximate}, where inference is also based on a discrepancy between empirical distributions and is therefore intrinsically permutation-invariant at the data level. In WABC, however, the parameter of interest is global, and the method is designed for samples viewed as unordered point clouds. As a consequence, the issue of recovering compartment-specific local parameters is not addressed there. By contrast, our setting is hierarchical: while permutation invariance can be exploited to improve acceptance, inference ultimately concerns not only the global component \( \globparam \) but also the local parameters \( \locparam_{1:K} \), at least up to relabelling.}

\new{From a computational viewpoint, \citet{bernton2019approximate} also discuss several approximations to exact optimal transport in the multivariate case, including entropically regularized Sinkhorn divergences, Hilbert-curve sorting, and a greedy swapping refinement initialized from the Hilbert ordering. Using their notation, their sample size \(n\) corresponds here to the number \(K\) of compartments, whereas their observation dimension \(d_y\) corresponds to the dimension of each compartment-level observation. These approximations are particularly attractive when the number of points is very large, but the Hilbert-based approximation is reported to be accurate mainly for small observation dimension and to deteriorate as \(d_y\) grows. In our applications, \(K\) is only moderately large---with a limiting case around \(K=100\), which already corresponds to the scale at which WABC considers the number of data points---while the dimension of each compartment-level observation may be non-negligible. In this regime, approximate transport schemes can lose accuracy quickly, and solving the assignment problem exactly remains a practical and robust option.}

However, this elegant result no longer holds in the multivariate case. When each $\obsdata_k$ is a vector in $\mathbb{R}^n$ with $n > 1$, there is no natural way to define a deterministic and permutation-optimal ordering. \old{As a result, we cannot rely on sorting to minimize $d(\obsdata, \simdata_{\sigma})$ in general, and alternative strategies are needed as soon as the dimension is at least 2.} Instead, we are compelled to solve the optimization problem directly, which leads us to the following permABC acceptance criterion:
\begin{equation}
    \label{eqopti}
    \min_{\sigma \in \mathcal{S}_K} d(\obsdata, \simdata_{\sigma}) \leq \varepsilon.
\end{equation}

This condition can be interpreted as accepting if our simulated dataset $\simdata$ matches, up to a permutation, the observed dataset $\obsdata$, and it effectively enhances the acceptance rate by a factor of $K!$ when $\varepsilon$ is small. This optimization task can be formulated as a two-dimensional linear assignment problem, or equivalently a bipartite matching problem, and is efficiently solved using a modified Jonker--Volgenant algorithm \citep{jonker1987shortest}, with a computational complexity of $\mathcal{O}(K^3)$ \citep{crouseImplementing2DRectangular2016a}. While this reduces the combinatorial cost from $K!$ to polynomial time, the $\mathcal{O}(K^3)$ complexity remains substantial. Nevertheless, it makes problems with up to $K = 100$ compartments tractable in practice (see the real-world application in Section~\ref{sec:sir}).

However, this approach does not yield samples from the pseudo-posterior defined in \eqref{eqabc}, but rather from the following permutation-invariant version:
\begin{equation} 
    \label{eqpermabc_tilde}
    \begin{split}
        \Tilde \pi_{\varepsilon} (\param \mid \obsdata)  & \propto \int_{\dataspace^K} \Tilde \pi_{\varepsilon} (\param , \simdata \mid \obsdata) \, d\simdata \\
        & \propto \pi(\param) \int_{\Tilde A_{\obsdata,\varepsilon}} f(\simdata \mid \param) \, d\simdata,
    \end{split}
\end{equation}
where the acceptance region is defined as
\[
\tilde{A}_{\obsdata,\varepsilon} = \left\{\simdata \in \dataspace^K \,\middle|\, \min_{\sigma \in \mathcal{S}_K} d(\obsdata, \simdata_{\sigma}) \leq \varepsilon \right\} = \bigcup_{\sigma \in \mathcal{S}_K} B_{\varepsilon}(\obsdata_{\sigma}),
\]
with \( B_{\varepsilon}(\obsdata_{\sigma}) \) denoting the \( \varepsilon \)-ball centered at \( \obsdata_{\sigma} \) for the distance \( d \).

In this formulation, the likelihood \( f(\obsdata \mid \param) \) is approximated by its integral over a neighborhood of all permutations of the observed dataset:
\[
\int_{\tilde{A}_{\obsdata,\varepsilon}} f(\simdata \mid \param) \, d\simdata \approx \sum_{\sigma \in \mathcal{S}_K} f(\obsdata_{\sigma} \mid \param).
\]

Inference in this setting targets the global parameter \( \globparam \), but is only defined up to a permutation of the local components \( \locparam_{1:K} \), effectively yielding posterior samples of the form \( \Tilde{\param} = (\globparam, \{\locparam_1, \dots, \locparam_K\})\), where the local parameters are treated as an unordered set. \old{This formulation is particularly suitable when the local parameters are considered nuisance variables, as in many hierarchical modeling contexts.} We discuss below how to recover a pseudo-posterior of each $\mu_k$ if needed.

To our knowledge, this permutation-based version of ABC has not been formally studied in the literature. We therefore introduce a novel methodology, termed \emph{permABC}, which builds upon this idea while addressing its limitations. In particular, our method enables proper inference over the full parameter vector \( \param \), including the ordering of local components, even in regimes where \( K \) is large, and does so while preserving a high acceptance rate.

\subsection{The permABC algorithm}
\label{secpermabc_star}


Restoring identifiability among compartments is essential for enabling inference on the full parameter vector. One natural way to achieve this is by applying a suitable permutation to the accepted parameters, such that the simulated data are optimally aligned with the observed dataset. However, explicitly testing all \( K! \) permutations is computationally intractable.

Fortunately, whenever the acceptance condition \eqref{eqopti} is satisfied, there exists at least one permutation of the simulated dataset that meets the ABC criterion, namely the optimal one minimizing the distance. This permutation, denoted by \( \sigma^* = \argmin_{\sigma \in \mathcal{S}_K} d(\obsdata, \simdata_{\sigma}) \), provides a natural way to reorder the parameter vector.

In practice, the \emph{permABC} algorithm proceeds similarly to standard ABC: we sample \( \parprop \sim \pi(\cdot) \) and simulate data \( \simdata \sim f(\cdot \mid \parprop) \). If the minimum distance over all permutations satisfies the tolerance condition, we accept and return the pair \( (\parprop_{\sigma^*}, \simdata_{\sigma^*}) \). This induces a new pseudo-posterior distribution, denoted \( \pi_{\varepsilon}^*(\cdot, \cdot \mid \obsdata) \), which corresponds to the pushforward of the distribution \( \tilde{\pi}_\varepsilon \) defined in \eqref{eqpermabc_tilde} under the optimal alignment map \( \proj \):

\begin{equation}
\label{eqproj}
\pi^{*}_{\varepsilon}(\cdot, \cdot \mid \obsdata)
:= \projsharp \tilde{\pi}_{\varepsilon}(\cdot, \cdot \mid \obsdata)
\end{equation}
where
\[
\proj: (\param, \simdata) \mapsto (\param_{\sigma^*}, \simdata_{\sigma^*}).
\]
By construction, \( \proj \circ \proj = \proj \), so \( \proj \) acts as a projection. Moreover, its preimage satisfies \( \proj^{-1}(\{\param, \simdata\}) = \{(\param_{\sigma}, \simdata_{\sigma}) \mid \sigma \in \mathcal{S}_K\} \).

As a result, the projected pseudo-posterior satisfies:
\begin{equation*}
    \begin{split}
        \pi_{\varepsilon}^* (\param, \simdata \mid \obsdata) &  \propto \Tilde \pi_{\varepsilon} (\proj^{-1}(\{(\param, \simdata)\}) \mid \obsdata)\cdot \mathbb I_{\text{Im}(\proj)}(\param,\simdata)\\
        & \propto \sum_{\sigma} \Tilde \pi_{\varepsilon} (\param_{\sigma}, \simdata_{\sigma} \mid \obsdata)\cdot \mathbb I_{\text{Im}(\proj)}(\param,\simdata)\\
        % & \propto \sum_{\sigma} \pi(\param_{\sigma}) f(\simdata_{\sigma} \mid \param_{\sigma}) \mathbb{I}_{\tilde A_{\obsdata,\varepsilon}}(\simdata_{\sigma})\cdot \mathbb I_{\text{Im}(\proj)}(\param,\simdata)\\
        % & \propto \pi(\param) f(\simdata \mid \param) \mathbb{I}_{\tilde A_{\obsdata,\varepsilon}}(\simdata)\cdot \mathbb I_{\text{Im}(\proj)}(\param,\simdata)\\
        & \propto \tilde \pi_{\varepsilon}(\param, \simdata\mid \obsdata) \cdot \mathbb I_{\text{Im}(\proj)}(\param,\simdata).
    \end{split}
\end{equation*}


\noindent which enforces a constraint that accepted samples must lie in the image of the projection, i.e., must be optimally aligned with the observed data.

This additional constraint effectively restricts the integration domain to a subset of datasets that are optimally ordered with respect to \( \obsdata \). We define this constrained space as
\[
\dataspace_{\obsdata}^{*K} = \left\{ \simdata \in \dataspace^K \mid d(\obsdata, \simdata) = \min_{\sigma \in \mathcal{S}_K} d(\obsdata, \simdata_{\sigma}) \right\}.
\]
Consequently, the marginal pseudo-posterior can be expressed as:
\begin{equation}
\label{eqpermabc_star}
\pi_{\varepsilon}^*(\param \mid \obsdata) \propto \pi(\param) \int_{A_{\obsdata,\varepsilon}^*} f(\simdata \mid \param) \, d\simdata,
\end{equation}
where the new acceptance region is 
\[
A_{\obsdata,\varepsilon}^* = B_{\varepsilon}(\obsdata) \cap \dataspace_{\obsdata}^{*K}.
\]

The \emph{permABC} algorithm (see Algorithm~\ref{alg:permABC_star}) therefore defines a different approximation of the posterior distribution than classical ABC methods. In particular, it yields a different posterior approximation  \( \pi_{\varepsilon}^*(\cdot \mid \obsdata) \), which we shall compare to the vanilla ABC approximation \( \pi_{\varepsilon}(\cdot \mid \obsdata) \).


\begin{algorithm}
\SetKwInOut{Inputt}{Input}
\SetKwInOut{Outputt}{Output}
\caption{permABC Vanilla}\label{alg:permABC_star}
\textbf{Input: }observed dataset $\obsdata$, number of iterations $N$, threshold $\varepsilon>0$.\\
\textbf{Output: }a sample $(\param^{(1)}, \dots, \param^{(N)})$ from the pseudo-posterior distribution $\pi_{\varepsilon}^*(\cdot\mid \obsdata)$.\\

\BlankLine

\For{$i = 1,\dots, N$}{
    \Repeat{$d(\obsdata, \simdata_{\sigma^*}) \leq \eps$}{
        $\parprop \sim \pi(\cdot)$\;
        $\simdata \sim f(\cdot \mid \parprop)$\;
        $\sigma^* = \argmin_{\sigma \in \mathcal{S}_K} d(\obsdata, \simdata_{\sigma})$\;
    }
    $\param^{(i)} = \param'_{\sigma^*}\;$
}
\end{algorithm}

\subsection{Comparison with ABC}
\label{sec:comparison}

The permABC method introduced in the previous section (Algorithm~\ref{alg:permABC_star}) differs from standard ABC in a key aspect: instead of accepting all permutations that satisfy the ABC criterion, it retains only the optimal one—namely, the permutation that minimizes the discrepancy defined in Equation~\eqref{eqopti}. This choice, motivated by computational efficiency, alters the pseudo-posterior distribution by constraining the pseudo-data to lie within $\dataspace^{*K}_{\obsdata}$, as previously described.


In contrast, a variant equivalent to standard ABC would consist in uniformly sampling a permutation among those satisfying the ABC acceptance condition and appropriately reweighting the contribution of each accepted sample by the size of this set. However, enumerating this set incurs a prohibitive $K!$ computational complexity. In Appendix~\ref{sec:appendix_strat}, we propose a stratified sampling strategy—termed \emph{permABC-strat}—to efficiently estimate this cardinality and recover the standard pseudo-posterior.

To clarify the connection between permABC and standard ABC, we establish a sufficient condition under which both methods coincide. Specifically, we prove that when the tolerance $\varepsilon$ is small enough, the set of permutations satisfying the ABC criterion contains at most one element. We formalize this in the following lemma:

\begin{lemma}
\label{lempermABC_star}
Let $\obsdata \in \dataspace^K$ denote the observed dataset, and define the discrepancy function $d(\cdot, \cdot)$ as in Equation~\eqref{eqdistance}. Let
\[
\varepsilon^* := \min_{\sigma \neq \mathrm{Id}} \frac{1}{2} d(\obsdata, \obsdata_\sigma).
\]
Then for any $\varepsilon < \varepsilon^*$ and any simulated dataset $\simdata \in \mathbb{R}^{K \times n}$, at most one permutation satisfies the ABC criterion:
\[
\mathrm{Card} \left( \left\{ \sigma \in \mathcal{S}_K : d(\obsdata, \simdata_\sigma) \leq \varepsilon \right\} \right) \leq 1.
\]
\end{lemma}

The proofs of the preceding lemma, the following theorem, and the subsequent corollary are deferred to Appendix~\ref{appendix_proofs}. 

We make the assumption that no pair of compartments have the exact same observed data. This assumption implies that $\varepsilon^*>0$, and will in particular be verified almost surely for continuous data. It could be relaxed to take into account ties, but we do not consider such an extension here. Under this assumption, we can show that permABC and standard ABC agree for small enough $\varepsilon$, as formalized in the following results.

\begin{theorem}
\label{th:permABC_star}
Let $\pi^*_{\varepsilon}( \cdot \mid \obsdata)$ be the pseudo-posterior distribution resulting from the permABC algorithm, and $\pi_{\varepsilon}( \cdot \mid \obsdata)$ that obtained via standard ABC. Then, there exists $\varepsilon^*>0$ such that for all $\varepsilon < \varepsilon^*$, we have:
\[
\pi^*_{\varepsilon}( \cdot \mid \obsdata) = \pi_{\varepsilon}( \cdot \mid \obsdata).
\]
\end{theorem}

This equivalence ensures that permABC retains the same asymptotic behavior as standard ABC in the limit of vanishing tolerance. We summarize this in the following corollary:

\begin{corollary}
\label{col:permabc}
If the pseudo-posterior distribution $\pi_{\varepsilon}(\cdot \mid \obsdata)$ obtained by the ABC algorithm converges to the true posterior distribution as $\varepsilon \to 0$, then the pseudo-posterior distribution $\pi_{\varepsilon}^*(\cdot \mid \obsdata)$ obtained by the permABC algorithm also converges to the true posterior distribution as $\varepsilon \to 0$.
\end{corollary}

Thus, the efficiency gains provided by permABC come at no cost to theoretical consistency. For foundational results on ABC convergence, we refer to \citet{fearnhead2018abc}, who provides general conditions under which ABC approximations converge to the true posterior, along with error bounds.




